\subsection*{Environment}

To streamline the configuration of the Croqtile build environment, developers can currently utilize the following command:

\begin{lstlisting}[language=croqtile,caption={Environment}]
make setup
\end{lstlisting}

This command retrieves the essential software prerequisites, such as the flex executable, bison executable (version 3.8 or higher), FileCheck utility, and gtest source code, among others. This process enables the building and testing of Croqtile.

Once the environment configuration is complete, Croqtile can be built using the command:

\begin{lstlisting}[language=croqtile,caption={Environment}]
make
\end{lstlisting}

Given that Croqtile is still in development, running unit tests is crucial to prevent using corrupted version:

\begin{lstlisting}[language=croqtile,caption={Environment}]
make test
\end{lstlisting}

In the event of any issues, the test should halt and report errors.

For development and testing purposes, the Makefile provides additional utilities:

\begin{lstlisting}[language=croqtile,caption={Environment}]
make help
\end{lstlisting}

displays available build and test targets. To test elementwise operators in the samples:

\begin{lstlisting}[language=croqtile,caption={Environment}]
make sample-test
\end{lstlisting}

runs all elementwise operator tests, while:

\begin{lstlisting}[language=croqtile,caption={Environment}]
make sample-test-operator OPERATOR=add
\end{lstlisting}

tests a specific operator (e.g., add, mul, relu, sigmoid, softmax, tanh).
